\section{Testovanie}
\label{chapter:testovanie}

V~tejto kapitole popisujeme proces testovania
implementovanej aplikácie a~uvádzame výsledky
overenia spoľahlivosti rozpoznávania jednotlivých
gest. Cieľom testovania bolo overiť, či hraničné
hodnoty stanovené v~kapitole~\ref{chapter:gesta}
zabezpečujú spoľahlivú detekciu gest pri ich
opakovanom vykonávaní v~rôznych podmienkach.

\subsection{Metodika testovania}

Testovanie sme rozdelili do troch fáz, ktoré sa
postupne približujú k~reálnym podmienkam použitia
aplikácie. V~prvej fáze sme overovali úspešnosť
rozpoznania jednotlivých gest pri ich izolovanom
vykonávaní. V~druhej fáze sme testovali kompletné
kompozitné sekvencie zodpovedajúce reálnym scenárom
použitia. V~tretej fáze sme overovali odolnosť
aplikácie voči falošne pozitívnym detekciám
pri bežnej aktivite ruky.

Testovanie prebiehalo v~debug móde Android aplikácie
(viď podsekciu~\ref{app:debug-final}), ktorý umožňuje
vizuálnu kontrolu rozpoznaných gest v~reálnom čase.
Pre každý test bola zaznamenávaná informácia o~tom,
či aplikácia gesto správne rozpoznala, prípadne aký
bol dôvod chybnej detekcie. Testy vykonával jeden
používateľ (autor práce) v~stoji, v~kontrolovanom
prostredí.

Pre vyhodnotenie sme použili nasledujúce metriky:
\begin{itemize}
  \item \textbf{Úspešnosť rozpoznania} -- podiel správne
        rozpoznaných gest z~celkového počtu vykonaných
        pokusov, vyjadrený v~percentách.
  \item \textbf{Falošne pozitívne detekcie} -- počet
        rozpoznaných gest v~situáciách, keď žiadne gesto
        nebolo úmyselne vykonané.
  \item \textbf{Úspešnosť kompozitných sekvencií} -- podiel
        sekvencií, pri ktorých výsledný stav aplikácie
        (skóre, režim) zodpovedal očakávaniu.
\end{itemize}

\subsection{Testovanie jednotlivých gest}
\label{subsec:test-jednotlive}

V~prvej fáze testovania sme každé gesto vykonali
20-krát v~ideálnych podmienkach (v~stoji, bez
sprievodných pohybov tela, v~neutrálnej polohe ruky
medzi pokusmi). Cieľom bolo overiť spoľahlivosť
detekcie samotného gesta bez vplyvu kontextu.

Pre testovanie gesta \uv{Hand Down} platí špecifický
postup -- toto gesto sa testuje vždy ako návrat
z~iného režimu, keďže v~režime \texttt{WAITING}
už ruka v~neutrálnej polohe je. Hand Down sme preto
testovali v~kombinácii s~aktivačným gestom Rise Arm.

Výsledky úspešnosti rozpoznania jednotlivých gest
sú zhrnuté v~tabuľke~\ref{tab:test-jednotlive}.

\begin{table}[H]
  \centering
  \caption{Úspešnosť rozpoznania jednotlivých gest}
  \label{tab:test-jednotlive}
  \renewcommand{\arraystretch}{1.3}
  \begin{tabularx}{\textwidth}{
      >{\raggedright\arraybackslash}X
      >{\centering\arraybackslash}p{0.15\textwidth}
      >{\centering\arraybackslash}p{0.18\textwidth}
      >{\centering\arraybackslash}p{0.18\textwidth}
    }
    \toprule
    \textbf{Gesto}         & \textbf{Pokusy} & \textbf{Rozpoznané} & \textbf{Úspešnosť} \\
    \midrule
    Rise Arm               & 20              & XX                  & XX,X\,\%           \\
    \midrule
    Hand Back              & 20              & XX                  & XX,X\,\%           \\
    \midrule
    Hand Down              & 20              & XX                  & XX,X\,\%           \\
    \midrule
    Flick                  & 20              & XX                  & XX,X\,\%           \\
    \midrule
    PASIVITA -- červený    & 20              & XX                  & XX,X\,\%           \\
    \midrule
    PASIVITA -- modrý      & 20              & XX                  & XX,X\,\%           \\
    \midrule
    NAPOMENUTIE -- červený & 20              & XX                  & XX,X\,\%           \\
    \midrule
    NAPOMENUTIE -- modrý   & 20              & XX                  & XX,X\,\%           \\
    \midrule
    TOUCHE                 & 20              & XX                  & XX,X\,\%           \\
    \midrule
    \textbf{Celkovo}       & \textbf{180}    & \textbf{XXX}        & \textbf{XX,X\,\%}  \\
    \bottomrule
  \end{tabularx}
\end{table}

% TODO: Doplň krátky komentár (3-5 viet) k výsledkom.
% Príklady čo môžeš spomenúť:
% - Ktoré gesto malo najvyššiu/najnižšiu úspešnosť a prečo
% - Aké boli typické príčiny chybných detekcií
% - Či výsledky zodpovedali očakávaniam stanoveným pri návrhu

Z~výsledkov vyplýva, že... % TODO

\subsection{Testovanie kompozitných sekvencií}
\label{subsec:test-sekvencie}

V~druhej fáze sme testovali kompletné scenáre
zodpovedajúce reálnemu použitiu aplikácie počas
zápasu. Každú sekvenciu sme vykonali 5-krát
a~overovali sme, či výsledný stav aplikácie
(skóre, aktívny režim, log udalostí) zodpovedá
očakávaniu. Sekvenciu sme považovali za úspešnú
iba ak prebehli všetky jej kroky správne -- ak
zlyhalo čo i len jedno gesto v~rámci sekvencie,
celá sekvencia sa hodnotila ako neúspešná.

Definované scenáre a~ich výsledky sú uvedené
v~tabuľke~\ref{tab:test-sekvencie}.

\begin{table}[H]
  \centering
  \caption{Úspešnosť rozpoznania kompozitných sekvencií}
  \label{tab:test-sekvencie}
  \renewcommand{\arraystretch}{1.3}
  \begin{tabularx}{\textwidth}{
      >{\raggedright\arraybackslash}X
      >{\centering\arraybackslash}p{0.13\textwidth}
      >{\centering\arraybackslash}p{0.15\textwidth}
    }
    \toprule
    \textbf{Sekvencia}                              & \textbf{Pokusy} & \textbf{Úspešnosť} \\
    \midrule
    Bodovanie červeného -- 1 bod                    & 5               & X/5                \\
    \midrule
    Bodovanie červeného -- 3 body                   & 5               & X/5                \\
    \midrule
    Bodovanie modrého -- 2 body                     & 5               & X/5                \\
    \midrule
    Napomenutie červeného + 2 body modrému          & 5               & X/5                \\
    \midrule
    Napomenutie modrého + 1 bod červenému           & 5               & X/5                \\
    \midrule
    Pasivita s timeoutom (automatický bod súperovi) & 5               & X/5                \\
    \midrule
    Pasivita zrušená bodovaním pred timeoutom       & 5               & X/5                \\
    \midrule
    TOUCHE                                          & 5               & X/5                \\
    \midrule
    \textbf{Celkovo}                                & \textbf{40}     & \textbf{XX/40}     \\
    \bottomrule
  \end{tabularx}
\end{table}

% TODO: Komentár k výsledkom kompozitných sekvencií.
% - Ktoré sekvencie boli najspoľahlivejšie?
% - Pri ktorých sa najčastejšie vyskytovali chyby?
% - Ako sa správal odpočet pasivity v scenároch F a G?

Pri testovaní sekvencií... % TODO

\subsection{Test falošne pozitívnych detekcií}
\label{subsec:test-false-positive}

V~tretej fáze sme testovali odolnosť aplikácie voči
nechceným detekciám. Počas 10-minútového obdobia
sme vykonávali bežné aktivity ruky bez úmyslu
vyvolať akékoľvek gesto. Tento test bol kľúčový pre
overenie použiteľnosti aplikácie počas reálneho
zápasu, kedy rozhodca robí množstvo iných pohybov
ruky (gestikulácia, manipulácia s~predmetmi, chôdza)
a~aplikácia by ich nemala interpretovať ako gestá.

Test prebiehal s~aplikáciou v~režime \texttt{WAITING}
a~zahŕňal nasledujúce aktivity:
\begin{itemize}
  \item chôdza po miestnosti (približne 2~minúty),
  \item sedenie a~písanie na klávesnici (približne 2~minúty),
  \item gestikulácia rukami pri rozprávaní (približne 2~minúty),
  \item manipulácia s~predmetmi -- pitie z~pohára,
        držanie telefónu (približne 2~minúty),
  \item bežné pohyby pri obliekaní a~uchopovaní predmetov
        (približne 2~minúty).
\end{itemize}

Počas testu aplikácia zaznamenala \textbf{X falošne
  pozitívnych detekcií}. % TODO: doplň počet a charakter
Charakter zaznamenaných detekcií je uvedený
v~tabuľke~\ref{tab:test-false-positive}.

\begin{table}[H]
  \centering
  \caption{Falošne pozitívne detekcie počas bežnej aktivity}
  \label{tab:test-false-positive}
  \renewcommand{\arraystretch}{1.3}
  \begin{tabularx}{\textwidth}{
      >{\raggedright\arraybackslash}p{0.30\textwidth}
      >{\centering\arraybackslash}p{0.15\textwidth}
      >{\raggedright\arraybackslash}X
    }
    \toprule
    \textbf{Rozpoznané gesto} & \textbf{Počet} & \textbf{Aktivita pri detekcii}    \\
    \midrule
    % TODO: Doplň riadky podľa skutočných výsledkov.
    % Príklad: Flick & 2 & Rýchle otočenie zápästia pri pití z pohára
    % Ak žiadne falošné detekcie nenastali, uveď to v texte
    % a tabuľku môžeš úplne vynechať.
    --                        & 0              & Žiadne falošne pozitívne detekcie \\
    \bottomrule
  \end{tabularx}
\end{table}

% TODO: Komentár k testu false positive.
% - Boli falošné detekcie len pri jednom type gesta?
% - Aký je dopad na praktické použitie?
% - Potvrdzuje to vhodnosť dvojfázového designu BODOVANIA?

Výsledky tretieho testu... % TODO

\subsection{Diskusia výsledkov}

% TODO: Toto je miesto pre tvoju interpretáciu výsledkov.
% Šablóna nižšie je orientačná - prepíš ju podľa toho,
% čo ti reálne vyjde.

Na základe výsledkov všetkých troch fáz testovania
môžeme konštatovať, že implementovaná aplikácia
spoľahlivo rozpoznáva navrhnuté gestá s~priemernou
úspešnosťou XX\,\%. % TODO: doplň priemer

Najvyššiu spoľahlivosť dosahujú gestá s~výraznou
a~unikátnou senzorovou signatúrou. % TODO: konkretizuj
% Príklad: Flick a TOUCHE -- prvé vďaka výraznému
% záporému špičkovému výkyvu na osi g_x, druhé
% vďaka kombinácii hodnôt na osiach a_x a a_y,
% ktorá sa neprekrýva so žiadnym iným gestom.

Naopak, mierne nižšiu spoľahlivosť sme zaznamenali
pri gestách... % TODO: konkretizuj a vysvetli prečo

Dvojfázový design gesta \uv{Bodovanie} (aktivačné
gesto + flick) sa potvrdil ako efektívny ochranný
mechanizmus -- v~teste falošne pozitívnych detekcií
neboli flicky nesprávne pripočítané v~žiadnom
prípade, kedy aplikácia nebola v~aktívnom bodovacom
režime. % TODO: prispôsob ak vyšlo inak

\subsection{Limitácie a~návrhy na zlepšenie}

Pri interpretácii výsledkov je potrebné zohľadniť
limitácie testovania a~implementácie:

\begin{itemize}
  \item \textbf{Kalibrácia na jednu osobu} -- hraničné
        hodnoty boli stanovené na základe meraní
        jedného používateľa. Pre rozhodcov s~odlišnou
        anatómiou ruky alebo individuálnym štýlom
        signalizácie môže byť potrebná individuálna
        kalibrácia.

  \item \textbf{Testovanie v~kontrolovaných podmienkach} --
        testy prebiehali v~stoji bez sprievodných
        pohybov celého tela. V~reálnom zápase sa
        rozhodca pohybuje po žinenke, čo môže ovplyvniť
        presnosť rozpoznávania.

  \item \textbf{Obmedzená vzorka testovania} -- vzhľadom
        na rozsah diplomovej práce nebolo možné vykonať
        testovanie počas reálneho zápasu s~viacerými
        rozhodcami. Validácia v~reálnych podmienkach
        zostáva otvorenou témou pre nadväzujúci výskum.

  \item \textbf{Absencia automatizovaného odpočtu pasivity
          v~produkčnom móde} -- ako bolo popísané
        v~kapitole~\ref{chapter:implementacia}, % TODO: over label
        produkčný mód nevyhodnocuje 30-sekundový
        odpočet automaticky z~dôvodu zložitosti
        validácie pasivity v~reálnom zápase.
\end{itemize}

Na základe identifikovaných limitácií navrhujeme
nasledujúce smery ďalšieho rozvoja:

\begin{itemize}
  \item Implementácia kalibračného režimu, v~ktorom
        si každý rozhodca môže prispôsobiť hraničné
        hodnoty individuálnej anatómii a~štýlu
        signalizácie.

  \item Nahradenie statických hraničných hodnôt
        modelom strojového učenia (napríklad neurónovou
        sieťou trénovanou na nameraných dátach), čo by
        mohlo zvýšiť robustnosť rozpoznávania pri rôznych
        používateľoch a~v~rôznych podmienkach.

  \item Rozsiahlejšie testovanie počas reálnych zápasov
        s~viacerými rozhodcami a~následná analýza
        zhody medzi rozhodnutiami aplikácie a~oficiálnym
        bodovým hodnotením.

  \item Integrácia s~oficiálnym bodovým systémom
        \Gls{uww}, čím by aplikácia mohla slúžiť
        ako podporný nástroj pre celý rozhodcovský tím.
\end{itemize}